\documentclass[11pt]{article}
% Shared preamble for the printable mini lessons (1.4 onward).
%   \documentclass[11pt]{article}
%   \input{../preamble}
% Build the key with:  pdflatex -jobname=<name>_KEY "\def\iskey{}\input{<name>}"
% Figures shared with the web page live in figures/Unit_X/<name>.tex and are
% pulled in with \fig{Unit_X/<name>}; the standalone package drops their own
% preamble, so their macros come from figures/preamble.tex, loaded here.
% Shared preamble for web figures.  Each figure is a standalone .tex that
% \input's this file; scripts/make-figures.sh turns it into an SVG.
\usepackage[T1]{fontenc}
\usepackage{helvet}
\renewcommand{\familydefault}{\sfdefault}
\usepackage{tikz}
\usetikzlibrary{decorations.pathreplacing}
% axis numbers in the text font, not math italic
\pgfkeys{/pgf/number format/.cd,assume math mode=true,1000 sep={}}

% dot plot: \dotplot{min}{max}{value/count,...}{axis label}
\newcommand{\dotplot}[4]{%
\begin{tikzpicture}[x=0.9cm,y=0.33cm]
  \draw[thick] (#1-0.5,0) -- (#2+0.5,0);
  \foreach \t in {#1,...,#2} {\draw (\t,0.25) -- (\t,-0.25) node[below] {\t};}
  \foreach \v/\c in {#3} {\ifnum\c>0 \foreach \k in {1,...,\c} {\fill (\v,\k*1.15) circle (0.16cm);}\fi}
  \node at ({(#1+#2)/2},-2.3) {#4};
\end{tikzpicture}}

% shape sketch histogram (12 bins): \shapehist{heights}{label}
\newcommand{\shapehist}[2]{%
  \foreach \h [count=\i from 0] in {#1} \draw[fill=gray!30] (\i,0) rectangle (\i+1,\h);
  \draw (-0.3,0) -- (12.3,0);
  \node[font=\bfseries] at (6,-2.6) {#2};}
\newcommand{\meanmark}[1]{\fill (#1,0) -- ++(-0.3,-0.75) -- ++(0.6,0) -- cycle;}
\newcommand{\medmark}[2]{\draw[very thick,dashed] (#1,0) -- (#1,#2);}

% number line with minor ticks: \nline{min}{max}{label step}{axis label}
\newcommand{\nline}[4]{%
  \draw[thick] (#1-0.5,0) -- (#2+0.5,0);
  \foreach \t in {#1,...,#2} \draw (\t,0.12) -- (\t,-0.12);
  \pgfmathtruncatemacro{\nxt}{#1+#3}
  \foreach \t in {#1,\nxt,...,#2} \draw (\t,0.2) -- (\t,-0.2) node[below] {\t};
  \node at ({(#1+#2)/2},-0.95) {#4};}
% box plot above the line: \bplot{min}{Q1}{med}{Q3}{max}
\newcommand{\bplot}[5]{%
  \draw[thick] (#1,1.0) -- (#2,1.0);
  \draw[thick] (#4,1.0) -- (#5,1.0);
  \draw[thick] (#1,0.75) -- (#1,1.25);
  \draw[thick] (#5,0.75) -- (#5,1.25);
  \draw[thick,fill=gray!20] (#2,0.55) rectangle (#4,1.45);
  \draw[very thick] (#3,0.55) -- (#3,1.45);}

% ---- added for 1.4 onward ----
% axis on y=0: \axis{min}{max}{label step}{tick step}{axis label}
\newcommand{\axis}[5]{%
  \draw[thick] (#1,0) -- (#2,0);
  \pgfmathsetmacro{\axT}{#1+#4}\foreach \t in {#1,\axT,...,#2} \draw (\t,0.1) -- (\t,-0.1);
  \pgfmathsetmacro{\axL}{#1+#3}\foreach \t in {#1,\axL,...,#2}
    \draw (\t,0.18) -- (\t,-0.18) node[below] {\pgfmathprintnumber{\t}};
  \node at ({(#1+#2)/2},-0.95) {#5};}
% box plot centred at height y: \bplotat{y}{min}{Q1}{med}{Q3}{max}
\newcommand{\bplotat}[6]{%
  \draw[thick] (#2,#1) -- (#3,#1);
  \draw[thick] (#5,#1) -- (#6,#1);
  \draw[thick] (#2,#1-0.22) -- (#2,#1+0.22);
  \draw[thick] (#6,#1-0.22) -- (#6,#1+0.22);
  \draw[thick,fill=gray!20] (#3,#1-0.42) rectangle (#5,#1+0.42);
  \draw[very thick] (#4,#1-0.42) -- (#4,#1+0.42);}
% an outlier dot: \outlier{x}{y}
\newcommand{\outlier}[2]{\fill (#1,#2) circle (2.4pt);}

\usepackage[margin=0.75in]{geometry}
\usepackage{array}
\usepackage{enumitem}
\usepackage{amsmath}
\usepackage[helvet]{sfmath}  % math in the same sans font as the text
\usepackage{standalone}
\pagestyle{empty}
\setlength{\parindent}{0pt}
\setlist[enumerate,1]{leftmargin=1.5em,itemsep=1em}
\setlist[enumerate,2]{label=(\alph*),leftmargin=2em,itemsep=2.2em}

\newif\ifkey
\ifdefined\iskey \keytrue \fi
\colorlet{keyred}{red!75!black}
\newcommand{\ans}[1]{\ifkey\leavevmode{\color{keyred}#1}\fi}
% vertical work space on the student copy, less on the key
\newcommand{\work}[1]{\ifkey\vspace{0.6em}\else\vspace{#1}\fi}
\newcommand{\fig}[1]{\input{../../figures/#1}}
% Title line: \lessontitle{Quartiles, IQR, and Box Plots}
\newcommand{\lessontitle}[1]{{\Large\bfseries #1\ifkey\ \textcolor{keyred}{KEY}\fi}\hfill Name \underline{\hspace{2in}}\par\vspace{0.6em}}
\newcommand{\partner}{\par\vspace{0.4em}Work with a partner. Compare answers when you finish each problem.\par}

% Unit quiz: two questions per mini lesson, each headed by the lesson's
% objective.  The key keeps every box the same height as the student copy,
% so key page N lines up with quiz page N.
\newcommand{\lessonhead}[1]{\par\vspace{0.6em}{\large\bfseries #1}\par\vspace{0.25em}\hrule\vspace{0.6em}}
\newcommand{\goal}[1]{\textbf{#1}\par\vspace{0.35em}}
% answer box: same height on both copies, answer printed only on the key
\newcommand{\abox}[2]{\par\begin{minipage}[t][#1][t]{\linewidth}\raggedright\vspace{0.2em}\ans{#2}\end{minipage}\par}
% circled choices: \opt{A} prints plain; \optans{B} is circled in red on the key
\newcommand{\opt}[1]{\tikz[baseline=(c.base)]\node[rounded corners=0.6em,inner sep=2.5pt](c){#1};}
% (the ring is an overlay, so it takes no space and the key keeps the student layout)
\newcommand{\optans}[1]{\ifkey\tikz[baseline=(c.base)]{\node[rounded corners=0.6em,inner sep=2.5pt,text=keyred](c){#1};\draw[overlay,keyred,thick,rounded corners=0.6em] (c.south west) rectangle (c.north east);}\else\opt{#1}\fi}
% centred fixed-width table column, so a blank cell is as wide as a filled one
\newcolumntype{C}[1]{>{\centering\arraybackslash}p{#1}}
\setlist[enumerate,2]{label=(\alph*),leftmargin=2em,itemsep=0.3em}
\setlist[enumerate,1]{leftmargin=1.5em,itemsep=1em}
\raggedright\hyphenpenalty=10000\exhyphenpenalty=10000

\begin{document}

\lessontitle{Unit 3 Quiz: Two-Way Tables}

Show your work.

% ======================================================================
\lessonhead{3.1 - Two-Way Tables and Relative Frequency}
\begin{enumerate}

\item \goal{I read a two-way table, and I fill in missing numbers in it.}
100 students in 9th and 10th grade chose their favorite sport. Each student chose one sport.
\begin{enumerate}
  \item Fill in the missing numbers. Show your work below the table.

  \begin{center}
  \renewcommand{\arraystretch}{2}
  \begin{tabular}{|l|C{0.95in}|C{0.95in}|C{0.95in}|C{0.95in}|}\hline
   & \textbf{Soccer} & \textbf{Basketball} & \textbf{Volleyball} & \textbf{Total} \\ \hline
  \textbf{9th grade} & 24 & \ans{17} & 11 & \ans{52} \\ \hline
  \textbf{10th grade} & \ans{21} & 15 & \ans{12} & 48 \\ \hline
  \textbf{Total} & 45 & \ans{32} & 23 & 100 \\ \hline
  \end{tabular}
  \end{center}
  \abox{11em}{10th grade, soccer: $45 - 24 = 21$ \qquad 10th grade, volleyball: $23 - 11 = 12$\\[0.4em]
  9th grade total: $100 - 48 = 52$ \qquad 9th grade, basketball: $52 - 24 - 11 = 17$\\[0.4em]
  Basketball total: $17 + 15 = 32$\\[0.4em]
  Check: $21 + 15 + 12 = 48$ \; and \; $45 + 32 + 23 = 100$}
  \item How many students are 10th graders who chose soccer?
  \abox{3.5em}{21 students}
  \item How many students did \textbf{not} choose volleyball?
  \abox{3.5em}{$100 - 23 = 77$ students}
  \item How many 9th graders did \textbf{not} choose soccer?
  \abox{3.5em}{$52 - 24 = 28$ students}
\end{enumerate}

\newpage
\item \goal{I move between a Venn diagram and a two-way table, and I make a relative frequency table.}
This Venn diagram shows the music that 40 students like.

\vspace{0.4em}
\begin{minipage}[t]{0.42\linewidth}\vspace{0pt}\raggedright
\scalebox{0.8}{\venn{Likes rap}{Likes pop}{12}{8}{14}{6}}
\end{minipage}\hfill
\begin{minipage}[t]{0.55\linewidth}\vspace{0pt}\raggedright
\begin{enumerate}
  \item Make a two-way table from the Venn diagram.

  \vspace{0.5em}
  \renewcommand{\arraystretch}{2}
  \begin{tabular}{|l|C{0.6in}|C{0.6in}|C{0.6in}|}\hline
   & \textbf{Pop} & \textbf{No pop} & \textbf{Total} \\ \hline
  \textbf{Rap} & \ans{8} & \ans{12} & \ans{20} \\ \hline
  \textbf{No rap} & \ans{14} & \ans{6} & \ans{20} \\ \hline
  \textbf{Total} & \ans{22} & \ans{18} & \ans{40} \\ \hline
  \end{tabular}
\end{enumerate}
\end{minipage}

\vspace{1em}
\begin{enumerate}[start=2,itemsep=0.9em]
  \item Make a relative frequency table. Divide each count by 40. Write percents.

  \begin{center}
  \renewcommand{\arraystretch}{2.6}
  \begin{tabular}{|l|C{1.15in}|C{1.15in}|C{1.15in}|}\hline
   & \textbf{Pop} & \textbf{No pop} & \textbf{Total} \\ \hline
  \textbf{Rap} & \ans{$\frac{8}{40} = 20\%$} & \ans{$\frac{12}{40} = 30\%$} & \ans{$\frac{20}{40} = 50\%$} \\ \hline
  \textbf{No rap} & \ans{$\frac{14}{40} = 35\%$} & \ans{$\frac{6}{40} = 15\%$} & \ans{$\frac{20}{40} = 50\%$} \\ \hline
  \textbf{Total} & \ans{$\frac{22}{40} = 55\%$} & \ans{$\frac{18}{40} = 45\%$} & \ans{100\%} \\ \hline
  \end{tabular}
  \end{center}
  \item What percent of all the students like at least one of the two kinds of music?
  \abox{5.5em}{$12 + 8 + 14 = 34$ students. \quad $\frac{34}{40} = 85\%$ \qquad (Or: $100\% - 15\% = 85\%$)}
  \item Which statement is supported by the relative frequency table? Circle the letter.

  \vspace{0.3em}
  \renewcommand{\arraystretch}{1.3}
  \begin{tabular}{@{}ll}
  \opt{A} & 8\% of the students like both rap and pop.\\
  \optans{B} & 55\% of the students like pop.\\
  \opt{C} & 30\% of the students like neither rap nor pop.\\
  \opt{D} & 35\% of the students like rap but not pop.\\
  \end{tabular}
\end{enumerate}

\end{enumerate}

% ======================================================================
\newpage
\lessonhead{3.2 - Marginal and Conditional Distributions}
\begin{enumerate}[resume]

\item \goal{I find a marginal distribution and a conditional distribution.}
80 students in 10th and 11th grade chose their favorite school subject.
\begin{center}
\renewcommand{\arraystretch}{1.6}
\begin{tabular}{|l|C{0.8in}|C{0.8in}|C{0.8in}|C{0.8in}|}\hline
 & \textbf{Math} & \textbf{Science} & \textbf{Art} & \textbf{Total} \\ \hline
\textbf{10th grade} & 12 & 9 & 9 & 30 \\ \hline
\textbf{11th grade} & 18 & 21 & 11 & 50 \\ \hline
\textbf{Total} & 30 & 30 & 20 & 80 \\ \hline
\end{tabular}
\end{center}
\begin{enumerate}
  \item Find the marginal distribution of favorite subject.
  \abox{9.5em}{Divide each total by the grand total, 80.\\[0.4em]
  Math $\frac{30}{80} = 37.5\%$ \qquad Science $\frac{30}{80} = 37.5\%$ \qquad Art $\frac{20}{80} = 25\%$}
  \item Find the conditional distribution of favorite subject for 11th graders.
  \abox{9.5em}{Divide by the 11th grade total, 50.\\[0.4em]
  Math $\frac{18}{50} = 36\%$ \qquad Science $\frac{21}{50} = 42\%$ \qquad Art $\frac{11}{50} = 22\%$}
  \item Is each statement a \textbf{marginal} or a \textbf{conditional} percent? Circle one.

  \vspace{0.6em}
  \renewcommand{\arraystretch}{1.8}
  \begin{tabular}{@{}p{3.6in}ll}
  25\% of all the students chose art. & \optans{Marginal} & \opt{Conditional}\\
  Of the 10th graders, 40\% chose math. & \opt{Marginal} & \optans{Conditional}\\
  \end{tabular}
\end{enumerate}

\newpage
\item \goal{I choose the right total to divide by, and I decide if two categorical variables are associated.}
A student recorded the weather and the school bus for 75 school days.
\begin{center}
\renewcommand{\arraystretch}{1.6}
\begin{tabular}{|l|C{0.9in}|C{0.9in}|C{0.9in}|}\hline
 & \textbf{Bus late} & \textbf{Bus on time} & \textbf{Total} \\ \hline
\textbf{Rain} & 9 & 6 & 15 \\ \hline
\textbf{No rain} & 12 & 48 & 60 \\ \hline
\textbf{Total} & 21 & 54 & 75 \\ \hline
\end{tabular}
\end{center}
\begin{enumerate}
  \item On what percent of the rainy days was the bus late?
  \abox{6.5em}{Divide by the rain total, 15. \quad $\frac{9}{15} = 60\%$}
  \item On what percent of the days with no rain was the bus late?
  \abox{6.5em}{Divide by the no rain total, 60. \quad $\frac{12}{60} = 20\%$}
  \item Is rain associated with a late bus? Explain.
  \abox{7.5em}{Yes. The percents are different. The bus was late on 60\% of the rainy days, but on only 20\% of the days with no rain.}
  \item Of the days when the bus was late, what percent had rain? Which fraction gives the answer? Circle the letter.

  \vspace{0.4em}
  \renewcommand{\arraystretch}{2.3}
  \begin{tabular}{@{}l@{\hspace{0.4em}}l@{\hspace{3em}}l@{\hspace{0.4em}}l@{\hspace{3em}}l@{\hspace{0.4em}}l@{\hspace{3em}}l@{\hspace{0.4em}}l}
  \opt{A} & $\dfrac{9}{15}$ & \optans{B} & $\dfrac{9}{21}$ & \opt{C} & $\dfrac{9}{75}$ & \opt{D} & $\dfrac{21}{75}$\\
  \end{tabular}
\end{enumerate}

\end{enumerate}

\end{document}
